% Part: sets-functions-relations
% Chapter: relations
% Section: reflections
%
\documentclass[../../../include/open-logic-section]{subfiles}

\begin{document}

\olfileid{sfr}{rel}{ref}
\olsection{Pamawas Filosofis}

Ing \olref[set]{sec}, kita nemtokake relasi minangka himpunan
tartamtu. Kita perlu mandheg sedhela lan ngajokake pitakon
filosofis: apa kang \emph{ditindakake} dening definisi kaya iki?
Banget anggone mangu-mangu yen kita arep nyatakake manawa kita
wis \emph{nemokake} kasunyatan identitas metafisis; tuladhane,
yen relasi tatanan ing $\Nat$ \emph{pranyata} iku himpunan
$R=  \Setabs{\tuple{n,m}}{n, m \in \Nat\text{ lan }
n < m}$ kang kita tetepake ing \olref[set]{sec}.
Ana telung alesan kanggo mangu-mangu iki.

Kapisan: ing \olref[set][pai]{wienerkuratowski}, kita netepake
$\tuple{a, b} =\{\{a\}, \{a, b\}\}$.
Saiki gatekna definisi alternatif
$\lVert a,b\rVert = \{\{b\}, \{a, b\}\} = \tuple{b,a}$.
Yen $a \neq b$, mula $\tuple{a, b} \neq \lVert a,b\rVert$.
Nanging kita uga bisa nggunakake $\lVert a,b\rVert$ minangka
definisi pasangan urut, tinimbang $\tuple{a,b}$.
Kaloro definisi padha-padha bisa digunakake kanthi becik.
Mula saiki ana rong calon kang padha apike kanggo "dadi"
relasi tatanan ing wilangan natural, yaiku:
\begin{align*}
		R &= \Setabs{\tuple{n,m}}{n, m \in \Nat \text{ lan }n < m}\\
		S &= \Setabs{\lVert n,m\rVert}{n, m \in \Nat \text{ lan }n < m}.
\end{align*}
Amarga $R \neq S$, miturut ekstensionalitas cetha yen ora
bisa \emph{loro-lorone} padha karo relasi tatanan ing~$\Nat$.
Nanging mung dadi pilihan sawenang-wenang, lan rada ngisin-isini,
yen ngaku manawa $R$, dudu $S$ (utawa kosok baline),
pancen \emph{iku} relasi tatanan minangka kasunyatan.
(Iki tuladha prasaja banget saka argumen nglawan reduksionisme
teori himpunan kang digawe misuwur dening Benacerraf ing
\citeyear{Benacerraf1965}. Kita bakal bali ngrembug kaping pirang-pirang.)

Kapindho: yen kita nganggep \emph{saben} relasi kudu
diidentifikasi karo himpunan, mula relasi kaanggotaan himpunan
dhewe, $\in$, kudune sawijining himpunan tartamtu.
Malah kudune himpunan $\Setabs{\tuple{x,y}}{x \in y}$.
Nanging apa himpunan iki ana? Kanthi ana Paradoks Russell,
pratelan yen himpunan kaya iki ana dudu prakara prasaja.
Nyatane, \oliflabeldef{cumul:::part}{teori himpunan kang kita
kembangake ing \olref[cumul][][]{part} bakal
\emph{nolak} anane himpunan iki.\footnote{Kanggo ndhisiki andharan,
iki alesane. Kanggo bukti kanthi kontradiksi, umpamakna
$I =\Setabs{\tuple{x,y}}{x \in y}$ ana.
Mula $\bigcup \bigcup I$ iku himpunan universal,
kang nalisir \olref[sfr][z][sep]{thm:NoUniversalSet}.}}{teori himpunan
bisa dikembangake kanthi tliti minangka teori aksiomatis,
lan teori kasebut pancen bakal nolak anane himpunan iki.}
Mula, sanajan sawetara relasi bisa dianggep himpunan,
relasi kaanggotaan himpunan kudu dadi kasus khusus.

Katelu: nalika "ngidentifikasi" relasi karo himpunan,
kita nyebut yen kita kena nulis $Rxy$ kanggo
$\tuple{x,y} \in R$. Iki trep, angger relasi kaanggotaan,
"$\in$", dianggep \emph{minangka} predikat.
Nanging yen "$\in$" dianggep nglambangake sawijining jinis
himpunan, ukara "$\tuple{x,y} \in R$" mung dumadi saka
telung istilah tunggal kang nglambangake himpunan:
"$\tuple{x,y}$", "$\in$", lan "$R$".
Dhaptar jeneng kaya iki ora luwih bisa ngandharake proposisi
tinimbang rerangken tanpa teges: "cangkir wadhah pulpen meja".
Maneh, sanajan sawetara relasi bisa dianggep himpunan,
relasi kaanggotaan himpunan kudu dadi kasus khusus.
(Iki nggabungake versi prasaja saka paradoks konsep
\emph{jaran} Frege lan sawijining bantahan misuwur
kang tau diajokake Wittgenstein marang Russell.)

Dadi, apa akibat saka kabeh iki? Ora ana kang \emph{luput}
yen kita nyebut relasi ing wilangan minangka himpunan.
Kita mung kudu ngerteni maksud nalika pratelan kasebut digawe.
Kita ora nyatakake kasunyatan identitas metafisis.
Kita mung nyathet yen, ing konteks tartamtu, kita bisa
(lan bakal) \emph{nganggep} relasi (tartamtu) minangka
himpunan tartamtu.

\end{document}
